\chapter{الاحتمالات والمعاينة}\label{ch:g10:proba}

يسند \omterm{def:g10:proba:distribution}{الاحتمال} أعدادًا إلى الصدفة. وانطلاقًا من تجارب ذات
إمكانيات متساوية \omterm{def:g10:proba:distribution}{الاحتمال} --- النرد، والقطع النقدية، والورق --- يبني هذا الفصل
القواعد الأساسية لحساب احتمالات الحوادث، ويقدّم
أشجار الاحتمالات من أجل التجارب ذات المرحلتين، وينتهي بالصلة بين
الاحتمالات والتواترات المرصودة. وتستمر النظرية في
\cref{ch:g11:prob}.

\section{التجارب العشوائية والحوادث}

\begin{definition}[الإمكانيات والحوادث]\label{def:g10:proba:events}
\emph{التجربة العشوائية}\index{تجربة عشوائية} تجربة لا يمكن
التنبؤ بنتيجتها (كرمي حجر نرد). ومجموعة كل النتائج الممكنة
هي \emph{مجموعة الإمكانيات}\index{مجموعة الإمكانيات}
$\Omega = \{\omega_1, \dots, \omega_n\}$؛ وعناصرها هي
\emph{الإمكانيات}. و\emph{الحادثة}\index{حادثة} مجموعة من الإمكانيات؛ و
\emph{تتحقق} الحادثة عندما تنتمي نتيجة التجربة إليها.
\end{definition}

\begin{definition}[التوزيع الاحتمالي]\label{def:g10:proba:distribution}
إسناد عدد $p_i \geq 0$ إلى كل إمكانية $\omega_i$، مع
$p_1 + \dots + p_n = 1$، يعرّف \emph{احتمالًا}\index{احتمال}:
فاحتمال $\P(A)$ \omterm{def:g10:proba:events}{الحادثة} $A$ هو مجموع الأعداد $p_i$ الموافقة
للإمكانيات الموجودة في $A$. وبوجه خاص $\P(\Omega) = 1$، \omterm{def:g10:proba:events}{والحادثة} المستحيلة
$\varnothing$ احتمالها $0$.
\end{definition}

\begin{proposition}[الإمكانيات المتساوية الاحتمال]\label{prop:g10:proba:uniform}
إذا كانت الإمكانيات $n$ متساوية \omterm{def:g10:proba:distribution}{الاحتمال} (أي $p_i = \frac1n$ لكل منها)، فإنه من أجل
كل \omterm{def:g10:proba:events}{حادثة} $A$:
\[
\P(A) = \frac{\text{عدد إمكانيات } A}{n}
= \frac{\text{الحالات الملائمة}}{\text{الحالات الممكنة}} .
\]
\end{proposition}

\begin{proof}
تحوي $A$ عددًا $k$ من الإمكانيات، \omterm{def:g10:proba:distribution}{احتمال} كل منها $\frac1n$،
إذن $\P(A) = \frac kn$.
\end{proof}

\begin{example}\label{ex:g10:proba:die}
ارمِ حجر نرد سليمًا: $\Omega = \{1, 2, 3, 4, 5, 6\}$، \omterm{def:g10:proba:distribution}{واحتمال} كل إمكانية
$\frac16$. \omterm{def:g10:proba:events}{والحادثة} $A$: ``النتيجة زوجية'' هي
$A = \{2, 4, 6\}$، إذن $\P(A) = \frac36 = \frac12$. \omterm{def:g10:proba:events}{والحادثة} $B$: ``النتيجة
لا تقل عن $5$'' هي $B = \{5, 6\}$، إذن $\P(B) = \frac26 =
\frac13$.
\end{example}

\section{تركيب الحوادث}

\begin{definition}[التقاطع والاتحاد والمتممة]\label{def:g10:proba:operations}
لتكن $A$ و $B$ حادثتين.
\begin{itemize}
  \item تتحقق $A \cap B$ (``$A$ و $B$'') عندما تتحقق الحادثتان معًا؛
  \item وتتحقق $A \cup B$ (``$A$ أو $B$'') عندما تتحقق إحداهما على الأقل؛
  \item و\emph{المتممة}\index{متممة} $\bar A$ (``ليس $A$'')
        تتحقق بالضبط عندما لا تتحقق $A$؛
  \item وتكون $A$ و $B$ \emph{متنافيتين}\index{حادثتان متنافيتان} إذا لم يمكن أن تتحققا
        معًا: $A \cap B = \varnothing$.
\end{itemize}
\end{definition}

\begin{theorem}[قاعدتا الجمع]\label{thm:g10:proba:addition}
من أجل كل حادثتين $A$ و $B$:
\[
\P(A \cup B) = \P(A) + \P(B) - \P(A \cap B),
\qquad
\P(\bar A) = 1 - \P(A) .
\]
وفي حالة حادثتين \omterm{def:g10:proba:operations}{متنافيتين}، تبسط القاعدة الأولى إلى
$\P(A \cup B) = \P(A) + \P(B)$.
\end{theorem}

\begin{proof}
يضيف مجموع $\P(A) + \P(B)$ \omterm{def:g10:proba:distribution}{احتمال} كل إمكانية من
$A \cup B$ مرة واحدة، \emph{عدا} إمكانيات $A \cap B$ التي تُحسب
مرتين --- مرة في $A$ ومرة في $B$. وطرح $\P(A \cap B)$ يصلح
العدّ المزدوج. أما \omterm{def:g10:proba:operations}{المتممة}: فالحادثتان $A$ و $\bar A$ متنافيتان
و $A \cup \bar A = \Omega$، إذن $\P(A) + \P(\bar A) = \P(\Omega) = 1$.
\end{proof}

\begin{omfigure}
\begin{tikzpicture}[scale=0.9]
  \draw[omExo, thick, rounded corners] (-3.1,-1.7) rectangle (3.4,1.7);
  \node[omExo, font=\small] at (3.0,1.45) {$\Omega$};
  \begin{scope}
    \clip (-0.8,0) circle (1.25);
    \fill[omProp!25] (0.8,0) circle (1.25);
  \end{scope}
  \draw[omDef, thick] (-0.8,0) circle (1.25);
  \draw[omThm, thick] (0.8,0) circle (1.25);
  \node[omDef, font=\small] at (-1.6,0.9) {$A$};
  \node[omThm, font=\small] at (1.7,0.9) {$B$};
  \node[font=\small] at (0,0) {$A \cap B$};
\end{tikzpicture}

{\small حادثتان متداخلتان: في $\P(A) + \P(B)$ يُحسب \omterm{def:g10:numbers:interunion}{التقاطع}
المظلل $A \cap B$ مرتين، وهذا يفسّر الطرح
في قاعدة الجمع.}
\end{omfigure}

\begin{example}\label{ex:g10:proba:cards}
اسحب ورقة واحدة من رزمة قياسية من $52$ ورقة. ولتكن $A$: ``الورقة
كوبة'' ($13$ ورقة) و $B$: ``الورقة ملك'' ($4$ أوراق). عندئذٍ
تكون $A \cap B$ هي ``ملك الكوبة'' (ورقة واحدة)، و
\[
\P(A \cup B) = \frac{13}{52} + \frac{4}{52} - \frac{1}{52} = \frac{16}{52}
= \frac{4}{13}.
\]
وكثيرًا ما تكون قاعدة \omterm{def:g10:proba:operations}{المتممة} أسرع طريق: \omterm{def:g10:proba:distribution}{فاحتمال} أن تكون الورقة
\emph{ليست} كوبة هو $1 - \frac{13}{52} = \frac34$.
\end{example}

\section{أشجار الاحتمالات}

\begin{method}[التجارب ذات المرحلتين]\label{met:g10:proba:tree}
من أجل تجربة تُنجز على مرحلتين:
\begin{enumerate}
  \item ارسم \emph{شجرة}\index{شجرة احتمالات}: فرع لكل نتيجة ممكنة
        للمرحلة الأولى، ثم من كل فرع منها فرع لكل
        نتيجة للمرحلة الثانية، واكتب \omterm{def:g10:proba:distribution}{الاحتمال} على كل فرع
        (ومجموع الفروع الخارجة من أي عقدة يساوي $1$)؛
  \item \omterm{def:g10:proba:distribution}{احتمال} \emph{مسار} (أي ورقة من أوراق \omterm{met:g10:proba:tree}{الشجرة}) هو
        جداء الاحتمالات على فروعه؛
  \item \omterm{def:g10:proba:distribution}{واحتمال} \omterm{def:g10:proba:events}{حادثة} هو مجموع احتمالات
        المسارات التي تحققها.
\end{enumerate}
\end{method}

\begin{example}\label{ex:g10:proba:tree}
يحوي كيس $3$ كرات حمراء و $2$ كرة زرقاء. اسحب كرة، ثم أعدها،
ثم اسحب من جديد. ويعطي كل سحب أحمر (R) \omterm{def:g10:proba:distribution}{باحتمال} $\frac35$ و
أزرق (B) \omterm{def:g10:proba:distribution}{باحتمال} $\frac25$.

\begin{omfigure}
\begin{tikzpicture}[scale=1.0]
  \coordinate (root) at (0,0);
  \coordinate (R) at (2.3,1.1);
  \coordinate (B) at (2.3,-1.1);
  \coordinate (RR) at (5.0,1.7);
  \coordinate (RB) at (5.0,0.5);
  \coordinate (BR) at (5.0,-0.5);
  \coordinate (BB) at (5.0,-1.7);
  \draw (root) -- (R) node[midway, above, font=\small, sloped] {$\frac35$};
  \draw (root) -- (B) node[midway, below, font=\small, sloped] {$\frac25$};
  \draw (R) -- (RR) node[midway, above, font=\small, sloped] {$\frac35$};
  \draw (R) -- (RB) node[midway, below, font=\small, sloped] {$\frac25$};
  \draw (B) -- (BR) node[midway, above, font=\small, sloped] {$\frac35$};
  \draw (B) -- (BB) node[midway, below, font=\small, sloped] {$\frac25$};
  \node[omThm, font=\small] at (R) [above left=-2pt] {R};
  \node[omDef, font=\small] at (B) [below left=-2pt] {B};
  \node[right, font=\small] at (RR)
    {RR:\ $\frac35 \times \frac35 = \frac{9}{25}$};
  \node[right, font=\small] at (RB)
    {RB:\ $\frac35 \times \frac25 = \frac{6}{25}$};
  \node[right, font=\small] at (BR)
    {BR:\ $\frac25 \times \frac35 = \frac{6}{25}$};
  \node[right, font=\small] at (BB)
    {BB:\ $\frac25 \times \frac25 = \frac{4}{25}$};
\end{tikzpicture}

{\small \omterm{met:g10:proba:tree}{شجرة} السحبين مع الإرجاع. ومجموع احتمالات المسارات
الأربعة يساوي $1$.}
\end{omfigure}

\omterm{def:g10:proba:distribution}{واحتمال} الحصول على كرتين من اللون نفسه هو
$\P(\text{RR}) + \P(\text{BB}) = \frac{9}{25} + \frac{4}{25} =
\frac{13}{25}$؛ \omterm{def:g10:proba:distribution}{واحتمال} الحصول على زرقاء واحدة على الأقل هو
$1 - \P(\text{RR}) = 1 - \frac{9}{25} = \frac{16}{25}$ (بقاعدة \omterm{def:g10:proba:operations}{المتممة}).
\end{example}

\begin{remark}
ودون إرجاع، تتعلق احتمالات السحب الثاني بالسحب
الأول: ففي الكيس أعلاه، بعد سحب كرة حمراء، لا تبقى سوى $2$ حمراء و $2$
زرقاء، فيحمل الفرع الثاني $\frac24$ و $\frac24$. وتعمل
طريقة \omterm{met:g10:proba:tree}{الشجرة} دون تغيير؛ ولا تتغير سوى الأعداد على فروع المستوى الثاني.
\end{remark}

\section{التواترات والمعاينة}

عندما تُكرَّر \omterm{def:g10:proba:events}{تجربة عشوائية} مرات كثيرة، يقترب \omterm{def:g10:stats:series}{التواتر} المرصود
\omterm{def:g10:proba:events}{لحادثة} من احتمالها --- ولهذا يكون \omterm{def:g10:proba:distribution}{الاحتمال}
مفيدًا في وصف العالم الواقعي.

\begin{definition}[العينة]\label{def:g10:proba:sample}
\emph{العينة}\index{عينة} ذات الحجم $n$ هي نتيجة تكرار
التجربة نفسها $n$ مرة على نحو مستقل. و\emph{\omterm{def:g10:stats:series}{التواتر} المرصود} \omterm{def:g10:proba:events}{لحادثة}
في العينة هو عدد مرات تحققها مقسومًا على $n$.
\end{definition}

\begin{example}\label{ex:g10:proba:sampling}
تعطي قطعة نقدية سليمة الصورة \omterm{def:g10:proba:distribution}{باحتمال} $0.5$. ورميها $100$ مرة
نادرًا ما يعطي $50$ صورة بالضبط: فالعينات \emph{تتذبذب}. وتعطي السلاسل
النموذجية من $100$ رمية تواترات مثل $0.46$ و $0.53$ و $0.49$ --- قريبة من
$0.5$ لكن غير مساوية لها. والعينات الأكبر تتذبذب أقل: فمع $10\,000$
رمية، تقع التواترات المرصودة عادةً على بعد نحو $0.01$ من $0.5$.
\end{example}

\begin{remark}[فترة التذبذب]
قاعدة تقريبية مفيدة، تُبرَّر في \cref{ch:g11:binom}: من أجل \omterm{def:g10:proba:sample}{عينة}
حجمها $n$ \omterm{def:g10:proba:events}{لحادثة} احتمالها $p$ (مع $n$ كبير بما يكفي)، يقع
\omterm{def:g10:stats:series}{التواتر} المرصود في \omterm{def:g10:numbers:interval}{الفترة}
\[
\intcc{p - \frac{1}{\sqrt n}}{\,p + \frac{1}{\sqrt n}}
\]
في نحو $95\%$ من العينات. وإذا وقع \omterm{def:g10:stats:series}{تواتر} مرصود بعيدًا خارج
هذه \omterm{def:g10:numbers:interval}{الفترة}، فشكّك في قيمة $p$: وهذه هي نقطة انطلاق
الاختبار الإحصائي.
\end{remark}

\begin{example}\label{ex:g10:proba:test}
يدّعي مصنع أن $10\%$ فقط من قطعه معيبة
($p = 0.1$). وفي \omterm{def:g10:proba:sample}{عينة} من $n = 400$ قطعة، $64$ قطعة معيبة: أي أن \omterm{def:g10:stats:series}{التواتر}
المرصود $\frac{64}{400} = 0.16$. \omterm{def:g10:numbers:interval}{وفترة} التذبذب هي
$\intcc{0.1 - \frac{1}{20}}{0.1 + \frac{1}{20}} = \intcc{0.05}{0.15}$،
والقيمة $0.16$ تقع خارجها: \omterm{def:g10:proba:sample}{فالعينة} تلقي شكًّا جديًّا على الادعاء.
\end{example}

\section{تمارين}

\begin{exercise}[$\star$]\label{exo:g10:proba:1}
يُرمى حجر نرد سليم. احسب \omterm{def:g10:proba:distribution}{احتمال} كل \omterm{def:g10:proba:events}{حادثة}: ``الحصول على
$6$''؛ ``الحصول على عدد فردي''؛ ``الحصول على $4$ على الأكثر''؛ ``الحصول على
$7$''.
\end{exercise}

\begin{exercise}[$\star$]\label{exo:g10:proba:2}
يحوي كيس $5$ كرات خضراء و $3$ صفراء و $2$ سوداء؛ وتُسحب كرة واحدة
عشوائيًا. احسب \omterm{def:g10:proba:distribution}{احتمال} أن تكون خضراء؛ وأن لا تكون سوداء؛ وأن تكون
خضراء أو صفراء.
\end{exercise}

\begin{exercise}[$\star$]\label{exo:g10:proba:3}
في فوج من $30$ تلميذًا، يدرس $18$ الإسبانية، ويدرس $10$ الألمانية، ويدرس
$4$ اللغتين معًا. يُختار تلميذ عشوائيًا. مستعملًا قاعدة الجمع،
احسب \omterm{def:g10:proba:distribution}{احتمال} أن يدرس هذا التلميذ إحدى اللغتين على الأقل،
ثم \omterm{def:g10:proba:distribution}{احتمال} أن لا يدرس أيًّا منهما.
\end{exercise}

\begin{exercise}[$\star$]\label{exo:g10:proba:4}
تُرمى قطعة نقدية سليمة مرتين. ارسم \omterm{met:g10:proba:tree}{شجرة} التجربة واحسب
\omterm{def:g10:proba:distribution}{احتمال} الحصول على صورتين؛ وعلى صورة واحدة بالضبط؛ وعلى صورة واحدة
على الأقل.
\end{exercise}

\begin{exercise}[$\star\star$]\label{exo:g10:proba:5}
يحوي كيس $4$ كرات حمراء و $6$ بيضاء. وتُسحب كرتان واحدة بعد
الأخرى \emph{دون} إرجاع. ارسم \omterm{met:g10:proba:tree}{الشجرة بالاحتمالات}
الصحيحة على المستوى الثاني، واحسب \omterm{def:g10:proba:distribution}{احتمال} سحب
كرتين حمراوين، ثم \omterm{def:g10:proba:distribution}{احتمال} سحب كرتين مختلفتي اللون.
\end{exercise}

\begin{exercise}[$\star\star$]\label{exo:g10:proba:6}
يُرمى حجرا نرد سليمان ويُسجَّل مجموعهما.
\begin{enumerate}
  \item فسّر لماذا يمكن أخذ \omterm{def:g10:proba:events}{مجموعة الإمكانيات} على أنها الأزواج $36$
        المتساوية \omterm{def:g10:proba:distribution}{الاحتمال} $(1,1), (1,2), \dots, (6,6)$.
  \item احسب \omterm{def:g10:proba:distribution}{احتمال} أن يكون المجموع $7$؛ وأن يكون $12$؛
        وأن يكون $4$ على الأكثر.
  \item ما المجموع الأكثر \omterm{def:g10:proba:distribution}{احتمالًا}؟
\end{enumerate}
\end{exercise}

\begin{exercise}[$\star\star$]\label{exo:g10:proba:7}
يعرض سؤال متعدد الخيارات $4$ إجابات، واحدة منها صحيحة. ويجيب تلميذ
عن سؤالين مستقلين من هذا النوع عشوائيًا تمامًا. احسب
\omterm{def:g10:proba:distribution}{احتمال} أن يجيب عنهما معًا إجابة صحيحة، وعن واحد بالضبط، ولا عن أي
منهما. وتحقق من أن مجموع الاحتمالات الثلاثة $1$.
\end{exercise}

\begin{exercise}[$\star\star$]\label{exo:g10:proba:8}
يدّعي سياسي أن نسبة الرضا عنه $50\%$. وفي استطلاع شمل $n = 900$ ناخبًا
اختيروا عشوائيًا، أبدى $396$ رضاهم.
\begin{enumerate}
  \item احسب \omterm{def:g10:stats:series}{التواتر} المرصود \omterm{def:g10:numbers:interval}{وفترة} التذبذب
        حول $p = 0.5$ من أجل $n = 900$.
  \item هل يلقي الاستطلاع شكًّا على الادعاء؟
\end{enumerate}
\end{exercise}

\begin{exercise}[$\star\star\star$]\label{exo:g10:proba:9}
تكلّف لعبة $2$ للمشاركة فيها: ترمي حجر نرد سليمًا وتربح $6$ إذا حصلت على
$6$، و $3$ إذا حصلت على $5$، ولا شيء فيما عدا ذلك.
\begin{enumerate}
  \item احسب \omterm{def:g10:proba:distribution}{احتمال} كل مبلغ مربوح ($6$ و $3$ و $0$).
  \item على \omterm{def:g10:stats:quartiles}{مدى} $600$ لعبة، كم تتوقع من كل نتيجة \omterm{def:g10:numbers:approx}{تقريبًا}؟
        احسب مجموع الربح المنتظر والخسارة المنتظرة: فهل تستحق اللعبة
        أن تُلعب؟
\end{enumerate}
\end{exercise}

\section{مسألة: الأبواب الثلاثة التي خدعت ألف
رياضي}

\begin{problem}[{مسألة نهاية الأسبوع --- مفارقة مونتي هول و
بنات عمها: كيف تغيّر المعلومة الاحتمال، وكيف تفصل العينات
بين الادعاءات والتمني}]
\label{pb:g10:proba:1}
سنة 1990 نشرت الكاتبة مارلين فوس سافانت أحجية عن
مقدّم برنامج مسابقات وثلاثة أبواب وسيارة وعنزتين ---
وأجابت عنها إجابة صحيحة. فكتب إليها عشرة آلاف قارئ، منهم مئات
من حملة الدكتوراه في الرياضيات، يصرّون على أنها مخطئة. تسلّحك هذه
المسألة بالأشجار (\cref{met:g10:proba:tree})،
وبالمتممات وقاعدتي الجمع
(\cref{thm:g10:proba:addition})، وتدعك تحسم الجدال
بنفسك، وتختم بالأداة التي تحكم في كل نزاع
كهذا عمليًّا: \omterm{def:g10:proba:sample}{العينة} المتذبذبة.

\medskip
\noindent\textbf{الجزء الأول --- القواعد، على سبيل التحمية.}
\begin{enumerate}
  \item ورقة واحدة من رزمة قياسية من $52$ ورقة: احسب
        $P(\text{ملك أو كوبة})$ بقاعدة الجمع ---
        وقل لماذا يكون $\frac{4}{52} + \frac{13}{52}$ وحده
        خطأً.
  \item ثلاث رميات لقطعة نقدية سليمة: احسب $P(\text{صورة واحدة على الأقل})$ \omterm{def:g10:proba:operations}{بالمتممة}.
  \item يحوي كيس $3$ كرات حمراء و $2$ زرقاء؛ وتُسحب كرتان
        دون إرجاع. \omterm{met:g10:proba:tree}{بشجرة}، احسب
        $P(\text{أحمران})$ و $P(\text{واحدة من كل لون})$.
  \item أكمل توزيع السؤال 3 بحساب
        $P(\text{أزرقان})$ وتحقق من أن مجموع الاحتمالات الثلاثة
        يساوي $1$.
  \item يبلّغ أحدهم، من أجل حادثتين \emph{\omterm{def:g10:proba:operations}{متنافيتين}}،
        أن $P(A) = 0.6$ و $P(B) = 0.5$. أدِنه بسطر
        واحد من الحساب.
\end{enumerate}

\medskip
\noindent\textbf{الجزء الثاني --- الأبواب الثلاثة.}
اللعبة: تختبئ سيارة خلف أحد ثلاثة أبواب، وتختبئ عنزتان خلف
الآخرين. تختار بابًا (وليكن الباب 1). فيفتح المقدّم --- وهو
\emph{يعلم} أين السيارة --- أحد البابين الآخرين، كاشفًا عنزةً
دائمًا، ثم يعرض عليك التبديل
إلى الباب المغلق الباقي.
\begin{enumerate}[resume]
  \item قبل حساب أي شيء: هل ينفع التبديل، أم يضر، أم
        لا يهم؟ اكتب جوابك الحدسي --- بصدق.
  \item والآن \omterm{met:g10:proba:tree}{الشجرة}، بالمناقشة على موضع السيارة الحقيقي
        (\omterm{def:g10:proba:distribution}{احتمال} كل حالة $\frac13$): إذا كانت السيارة خلف بابك،
        فالتبديل يخسر؛ وإذا كانت خلف أي من البابين الآخرين، فماذا
        يجب أن يفعل المقدّم، وماذا يربح التبديل عندئذٍ؟
        احسب $P(\text{الفوز بالتبديل})$.
  \item الصيغة ذات السطر الواحد: يربح التبديل بالضبط عندما يكون
        اختيارك الأول خاطئًا. استنتج من جديد.
  \item والآن المتغيّر الحاسم --- ``مونتي الساقط'': لا
        \emph{يعلم} المقدّم أين السيارة، ويفتح أحد
        البابين الآخرين عشوائيًا، فتكشف الصدفة عنزةً. اذكر
        السيناريوهات الستة المتساوية \omterm{def:g10:proba:distribution}{الاحتمال} (موضع السيارة $\times$
        الباب المفتوح)، واستبعد تلك التي تُكشف فيها السيارة،
        واحسب \omterm{def:g10:proba:distribution}{احتمال} أن يربح التبديل
        \emph{بين السيناريوهات الباقية}. ما الذي تغيّر، و
        لماذا تهم معرفة المقدّم؟
  \item مضخة الحدس: مئة باب، تختار واحدًا، ويفتح
        المقدّم العليم $98$ بابًا خلفها عنزات. أتبدّل؟ وبأي
        \omterm{def:g10:proba:distribution}{احتمال} للفوز؟
  \item تواجه متشككًا يصرّ، مثل كتّاب الرسائل الألف،
        على ``النصف بالنصف''. صف محاكاة بورق اللعب
        (ثلاث أوراق، جولات كثيرة، صديق يقوم بدور المقدّم) وقل
        ما نسبة الفوز بالتبديل التي تتوقعها بعد $300$
        جولة --- ومعك \omterm{def:g10:numbers:interval}{فترة} تذبذب الفصل
        هامشًا.
\end{enumerate}

\medskip
\noindent\textbf{الجزء الثالث --- بنات عم المفارقة.}
\begin{enumerate}[resume]
  \item لأسرة طفلان؛ وتعلم أن \emph{أحدهما على الأقل
        ولد}. اذكر تشكيلات الطفلين المتساوية \omterm{def:g10:proba:distribution}{الاحتمال}
        الملائمة للمعلومة، واحسب
        $P(\text{ولدان})$.
  \item الأسرة نفسها، لكن المعلومة الآن: \emph{الطفل
        الأكبر ولد}. أعد الحساب. لماذا يختلف الجوابان،
        وما الذي يشتركان فيه مع صيغتَي مونتي؟
  \item ثلاث قطع نقدية سليمة: احسب $P(\text{الثلاثة متماثلة})$.
        وقد جادل الرياضي دالمبير سنة 1754 لصالح جواب آخر:
        ``إما أن تكون كلها متماثلة وإما لا ---
        حالة من حالتين.'' سمِّ الخطأ بالضبط، ومعك
        \cref{prop:g10:proba:uniform}.
  \item تكلّف لعبة $2$ يورو: ترمي حجر نرد وتربح $6$ يورو على
        الستة، و $3$ يورو على الخمسة، ولا شيء فيما عدا ذلك
        (\cref{exo:g10:proba:9}). على \omterm{def:g10:stats:quartiles}{مدى} $600$ لعبة، أي
        ربح \omterm{def:g10:stats:mean}{متوسط} لكل لعبة يتنبأ به العدّ؟ ويلعب
        المؤمِّن الرياضيات نفسها: \omterm{def:g10:proba:distribution}{احتمال} سنوي قدره $1\,\%$ لمطالبة
        بقيمة $10\,000$ يورو --- فما القسط
        \emph{العادل}، ولماذا يتجاوزه القسط الحقيقي؟
\end{enumerate}

\medskip
\noindent\textbf{الجزء الرابع --- العينات حكمًا.}
\begin{enumerate}[resume]
  \item تُرمى قطعة نقدية $100$ مرة. مستعملًا \omterm{def:g10:numbers:interval}{فترة} التذبذب
        $p \pm \frac{1}{\sqrt n}$، قرّر أي هذه
        النتائج ينبغي أن يثير الريبة: $55$ صورة؛ أم $70$
        صورة.
  \item يعطي استطلاع شمل $1\,000$ ناخب مرشحًا $52\,\%$.
        احسب \omterm{def:g10:numbers:interval}{فترة} التذبذب حول $50\,\%$ وقل
        هل \emph{يبرهن} الاستطلاع على أن المرشح
        متقدم.
  \item يدّعي مصنع أن $4\,\%$ من قطعه معيبة. وتُظهر \omterm{def:g10:proba:sample}{عينة} من
        $400$ قطعة $28$ قطعة معيبة ($7\,\%$)؛ وتُظهر \omterm{def:g10:proba:sample}{عينة} لاحقة من
        $2\,500$ قطعة $175$ قطعة ($7\,\%$ من جديد). اختبر الادعاء
        إزاء كل \omterm{def:g10:proba:sample}{عينة}. ولماذا يختلف الحكمان؟
  \item يلاحق كل استطلاع مصدرا خطأ: \emph{تحيّز}
        المعاينة (مسألة نهاية الأسبوع عن كيف تكذب المعطيات في الكتاب السابق) و\emph{تذبذب}
        المعاينة (هذا الفصل). أيهما ينكمش كلما نما
        $n$، وأيهما يصمد أمام أي حجم \omterm{def:g10:proba:sample}{عينة}؟ جملة
        واحدة لكل منهما.
  \item الخاتمة --- قائمة الاحتمالي، سطر لكل
        بند: تحقق من تساوي \omterm{def:g10:proba:distribution}{احتمال} الإمكانيات قبل العدّ
        (السؤال 15)؛ والأشجار والمتممات من أجل التجارب متعددة المراحل و
        ``على الأقل'' (السؤالان 2--3)؛ و\emph{المعلومة
        تغيّر الاحتمالات، وطريقة الحصول عليها تهم}
        (الأسئلة 7--14)؛ والعينات تتذبذب بمقدار نحو
        $\frac{1}{\sqrt n}$، والمحاكاة تحكم
        (الأسئلة 11 و 16--18).
\end{enumerate}
\end{problem}
